\documentclass[11pt,a4paper]{article}

\usepackage[utf8]{inputenc}
\usepackage[indonesian]{babel}
\usepackage[a4paper,top=2.2cm,bottom=2.2cm,left=2.4cm,right=2.4cm]{geometry}
\usepackage{graphicx}
\usepackage{array}
\usepackage{enumitem}
\usepackage{xcolor}
\usepackage{setspace}
\usepackage[breaklinks=true,colorlinks=true,linkcolor=primaryblue,urlcolor=primaryblue]{hyperref}

\definecolor{primaryblue}{RGB}{24, 43, 73}
\definecolor{secondaryblue}{RGB}{41, 128, 185}
\onehalfspacing

\begin{document}
\thispagestyle{empty}

\begin{center}
    \includegraphics[width=2.0cm]{../assets/logo_unib.png}\\[0.25cm]
    {\large \textbf{\color{primaryblue}SURAT PERNYATAAN KEASLIAN CIPTAAN}}\\[0.08cm]
    {\normalsize \textbf{\color{secondaryblue}PENGAJUAN PENCATATAN HAK CIPTA (DJKI / LPPM UNIVERSITAS BENGKULU)}}
\end{center}

\vspace{0.3cm}
\noindent Yang bertanda tangan di bawah ini:

\vspace{0.15cm}
\noindent
\begin{tabular}{p{4.3cm}p{0.4cm}p{10.8cm}}
Nama Lengkap & : & \textbf{Ir. Novalio Daratha, S.T., M.Sc., Ph.D.} \\
NIP / NIDN & : & 197811102003121002 \\
Jabatan / Unit Kerja & : & Dosen Program Studi S1 Teknik Elektro, Fakultas Teknik, Universitas Bengkulu \\
Alamat Instansi & : & Jl. W.R. Supratman, Kandang Limun, Kota Bengkulu 38371, Indonesia \\
Kewarganegaraan & : & Indonesia \\
\end{tabular}

\vspace{0.3cm}
\noindent Dengan ini menyatakan bahwa karya ciptaan yang saya mohonkan pencatatan Hak Ciptanya melalui Lembaga Penelitian dan Pengabdian kepada Masyarakat (LPPM) Universitas Bengkulu, yaitu:

\vspace{0.15cm}
\begin{enumerate}[label=\textbf{\arabic*.}, leftmargin=1.8em, itemsep=4pt]
    \item \textbf{Program Komputer:} \textit{Portal Pembelajaran Interaktif \& Virtual Forecasting Studio Analisis Deret Waktu dan Peramalan Sistem Tenaga (TEE-412)} \\
    (URL Publik: \url{https://www.ndaratha.my.id/deret-waktu} \;|\; Kode Sumber: \url{https://github.com/novatha/deret-waktu-tee412})
    \item \textbf{Buku / Modul Ajar:} \textit{Ekosistem Bahan Ajar, Buku Praktikum 12 Bab, Panduan Proyek Akhir IEEE, \& Bank Soal OBE (C1--C6) Mata Kuliah Analisis Deret Waktu dan Peramalan (TEE-412)}
\end{enumerate}

\vspace{0.15cm}
\noindent adalah benar-benar merupakan karya ciptaan asli hasil pemikiran dan penyusunan saya sendiri, dengan ketentuan sebagai berikut:
\begin{enumerate}[label=\textbf{\alph*.}, leftmargin=1.8em, itemsep=3pt]
    \item Ciptaan tersebut belum pernah dan tidak sedang menjadi objek sengketa kepemilikan hak kekayaan intelektual dengan pihak mana pun.
    \item Ciptaan tersebut tidak meniru, menjiplak, atau melanggar hak cipta pihak lain, serta seluruh pustaka/perangkat lunak sumber terbuka yang dirujuk telah dicantumkan secara patut sesuai etika akademik.
    \item Apabila di kemudian hari terbukti pernyataan ini tidak benar, saya bersedia bertanggung jawab penuh secara hukum dan melepaskan DJKI serta Universitas Bengkulu dari segala tuntutan pihak ketiga.
\end{enumerate}

\vspace{0.25cm}
\noindent Demikian surat pernyataan ini dibuat dengan sebenarnya dalam keadaan sadar dan tanpa paksaan untuk digunakan sebagai persyaratan pencatatan Hak Cipta di Direktorat Jenderal Kekayaan Intelektual (DJKI), Kementerian Hukum Republik Indonesia.

\vspace{0.8cm}
\noindent
\begin{tabular}{p{8.2cm}p{7.3cm}}
Mengetahui, & Bengkulu, November 2026 \\
Ketua Jurusan Teknik Elektro FT UNIB & Yang Menyatakan (Pencipta), \\[0.45cm]
 & {\footnotesize \color{gray}[ Materai Rp 10.000 ]} \\[1.25cm]
\textbf{\underline{Ir. Afriyastuti Herawati, S.T., M.T.}} & \textbf{\underline{Ir. Novalio Daratha, S.T., M.Sc., Ph.D.}} \\
NIP. 197404162000032001 & NIP. 197811102003121002 \\
\end{tabular}

\end{document}
